\section*{Limitations}

While \ours establishes a strong foundation for unified multimodal retrieval, we identify several areas for future improvement. 
(1) \emph{Language and Cultural Bias:} The training corpus utilized in this work is predominantly skewed toward English and Chinese. Consequently, the sparse head activations exhibit limited cross-lingual generalization. Extending the framework's sparse capabilities to a broader range of languages will necessitate highly language-diverse, large-scale training data. 
(2) \emph{Vocabulary Stability and Artifacts:} Unlike traditional masked language models with constrained vocabularies, operating over the expansive vocabulary of modern LLMs occasionally leads to anomalous token activations (\eg non-standard subwords or artifacts like ``\_alt''). Future iterations can explore targeted vocabulary pruning, refined semantic clustering, or post-hoc filtering mechanisms to ensure representational stability in strict production environments. 
(3) \emph{Modality-Specific Performance Gaps:} While the sparse and dense representations maintain close performance parity on text and static visual documents, we observe a slightly more pronounced gap in the video domain. We attribute this to the high information density and temporal dynamics inherent to video frames. This suggests that naively pooling spatiotemporal data into a flat sparse vector might encounter capacity bottlenecks. 